% \iffalse
%
% Copyright 2016-2025, Association for Computing Machinery
% This work may be distributed and/or modified under the
% conditions of the LaTeX Project Public License, either
% version 1.3 of this license or (at your option) any
% later version.
% The latest version of the license is in
%    http://www.latex-project.org/lppl.txt
% and version 1.3 or later is part of all distributions of
% LaTeX version 2005/12/01 or later.
%
% This work has the LPPL maintenance status `maintained'.
%
% The Current Maintainer of this work is Boris Veytsman,
% <borisv@lk.net>
%
% This work consists of the file acmart.dtx, the derived file
% acmart.cls, the files ACM-Reference-Format.bst, and templates
% sample-acmlarge.tex, sample-acmsmall.tex, sample-acmtog.tex,
% samplebody-conf.tex, samplebody-journals.tex, sample-manuscript.tex,
% sample-sigconf-authordraft.tex, sample-sigconf.tex,
% sample-sigplan.tex
%
% \fi
%
%
%% \CharacterTable
%%  {Upper-case    \A\B\C\D\E\F\G\H\I\J\K\L\M\N\O\P\Q\R\S\T\U\V\W\X\Y\Z
%%   Lower-case    \a\b\c\d\e\f\g\h\i\j\k\l\m\n\o\p\q\r\s\t\u\v\w\x\y\z
%%   Digits        \0\1\2\3\4\5\6\7\8\9
%%   Exclamation   \!     Double quote  \"     Hash (number) \#
%%   Dollar        \$     Percent       \%     Ampersand     \&
%%   Acute accent  \'     Left paren    \(     Right paren   \)
%%   Asterisk      \*     Plus          \+     Comma         \,
%%   Minus         \-     Point         \.     Solidus       \/
%%   Colon         \:     Semicolon     \;     Less than     \<
%%   Equals        \=     Greater than  \>     Question mark \?
%%   Commercial at \@     Left bracket  \[     Backslash     \\
%%   Right bracket \]     Circumflex    \^     Underscore    \_
%%   Grave accent  \`     Left brace    \{     Vertical bar  \|
%%   Right brace   \}     Tilde         \~}
%
%
% \MakeShortVerb{|}
% \def\guide{acmguide}
% \iffalse
% From
% http://tex.stackexchange.com/questions/117892/can-i-convert-a-string-to-catcode-11 by egreg
% \fi
% \begingroup
%  \everyeof{\noexpand}
%  \endlinechar=-1
%  \xdef\currentjob{\scantokens\expandafter{\jobname}}
% \endgroup
%
% \ifx\currentjob\guide\OnlyDescription\fi
% \GetFileInfo{acmart.dtx}
% \title{\LaTeX{} Class for the \emph{Association for Computing
% Machinery}\thanks{\copyright 2016--2025, Association for Computing Machinery}}
% \author{Boris Veytsman\thanks{%
% \href{mailto:borisv@lk.net}{\texttt{borisv@lk.net}},
% \href{mailto:boris@varphi.com}{\texttt{boris@varphi.com}}}}
% \date{\filedate, \fileversion}
% \maketitle
% \begin{abstract}
%   This package provides a class for typesetting publications of
%   the Association for Computing Machinery.
% \end{abstract}
% \tableofcontents
%
% \clearpage
%
%\section{Introduction}
%\label{sec:intro}
%
% The Association for Computing
% Machinery\footnote{\url{http://www.acm.org/}} is the world's largest
% educational and scientific computing society, which delivers
% resources that advance computing as a science and a
% profession.  It was one of the
% early adopters of \TeX\ for its typesetting.
%
% It provided several different classes for a number of journals and
% conference proceedings.  Unfortunately during the years since these
% classes were written, the code was patched many times, and
% supporting different versions of the classes became difficult.
%
% This new consolidated template package replaces all previous
% independent class files and packages and provides a single
% up-to-date LaTeX package with optional calls. The package uses only
% free \TeX\ packages and fonts included in \TeX Live, Mik\TeX\ and
% other popular \TeX\ distributions. The new ACM templates use a new
% font set (libertine) which will need to be installed on your machine
% before using the templates. Please download and install the
% libertine font set before writing your paper. Fonts used in the
% template cannot be substituted; margin adjustments are not allowed. 
%
%  I am grateful to
%  Michael D.~Adams,
%  Leif Andersen,
%  Lawrence Christopher Angrave,
%  Dirk Beyer,
%  Andrew Black,
%  Joachim Breitner,
%  Yegor Bugayenko,
%  Benjamin Byholm,
%  John Collins,
%  Roberto Di Cosmo,
%  Nils Anders Danielsson,
%  Michael Ekstrand,
%  Matthew Fluet,
%  Paolo G.~Giarrusso,
%  Ben Greenman,
%  Enrico Gregorio,
%  Jamie Davis,
%  Ulrike Fischer,
%  Jason Hemann,
%  Peter Kemp,
%  Luis Leiva,
%  Ben Liblit,
%  Rholais Lii,
%  LianTze Lim,
%  Kuldeep S. Meel,
%  Kai Mindermann,
%  Frank Mittelbach,
%  Serguei Mokhov,
%  Ross Moore,
%  John Owens,
%  Joel Nider,
%  Scott Pakin,
%  Tobias Pape,
%  Henning Pohl,
%  Philip Quinn,
%  Mathias Rav,
%  Andreas Reichinger,
%  Matteo Riondato,
%  Craig Rodkin,
%  Bernard Rous,
%  Feras Saad,
%  Kerry A. Seitz, Jr.,
%  M Senthilkumar,
%  David Shamma,
%  Gabriel Scherer,
%  Kartik Singhal,
%  Christoph Sommer,
%  Stephen Spencer,
%  Shin Hwei Tan,
%  Daniel Thomas,
%  Shari Trewin,
%  Zack Weinberg,
%  John Wickerson
%  and many others for their invaluable help.
%
% The development version of the package is available at
% \url{https://github.com/borisveytsman/acmart}.
%
%\section{User's guide}
%\label{sec:ug}
%
%
% This class uses many commands and customizaton options, so it might
% appear intimidating for a casual user.  Do not panic!  Many of these
% commands and options can be safely left with their default values
% or the values recommended by your conference or journal editors.  If
% you have problems or questions, do not hesitate to ask me directly
% or the community at \url{https://github.com/borisveytsman/acmart},
% \url{https://tex.stackexchange.com} or the closest \TeX\ Users
% Group.  The world-wide \TeX\ Users Group is at
% \url{https://tug.org/}; please consider joining us if you use \TeX\
% regularly.
%
%\subsection{Installation}
%\label{sec:ug_install}
%
% Most probably, you already have this package installed in your
% favorite \TeX\ distribution;  if not, you may want to upgrade.  You
% may need to upgrade it anyway since this package uses a number of
% relatively recent packages, especially the ones related to fonts.
%
% The latest released version of this package can be found on CTAN:
% \url{https://www.ctan.org/pkg/acmart}.   The development version can
% be found on GitHub: \url{https://github.com/borisveytsman/acmart}.
% At this address you can file a bug report---or even contribute your
% own enhancement by making a pull request.
%
% Please note that the version on Github is a development (or
% experimental) version: please download it for testing new features.
% The production version is the one on CTAN and ACM sites.
%
% Most users should not attempt to install this package themselves
% but should rather rely on their \TeX\ distributions to provide it.  If you
% decide to install the package yourself, follow the standard rules:
% \begin{enumerate}
% \item Run |latex acmart.ins|.  This will produce the file
% |acmart.cls|
% \item Put the files |acmart.cls|, |acm-jdslogo.png|,
%   and |ACM-Reference-Format.bst|
%   in places where \LaTeX{} can find them (see \cite{TeXFAQ} or
%   the documentation for your \TeX{} system).\label{item:install}
% \item Update the database of file names.  Again, see \cite{TeXFAQ}
% or the documentation for your \TeX{} system for the system-specific
% details.\label{item:update}
% \item The file |acmart.pdf| provides the documentation for the
% package.  (This is probably the file you are reading now.)
% \end{enumerate}
% As an alternative to items~\ref{item:install} and~\ref{item:update}
% you can just put the files in the working directory where your
% |.tex| file is.
%
%
% This class uses a number of other packages.  They are included in all
% major \TeX\ distributions (\TeX Live, Mac\TeX, Mik\TeX) of 2015 and
% later, so you probably have them installed.  Just in case here is
% the list of these packages:
% \begin{itemize}
% \item \textsl{amsart}, \url{https://ctan.org/pkg/amsart},
% \item \textsl{babel}, \url{https://ctan.org/pkg/babel},
% \item \textsl{balance}, \url{https://ctan.org/pkg/balance},
% \item \textsl{booktabs}, \url{https://ctan.org/pkg/booktabs},
% \item \textsl{caption}, \url{https://ctan.org/pkg/caption},
% \item \textsl{cmap}, \url{https://ctan.org/pkg/cmap},
% \item \textsl{comment}, \url{https://ctan.org/pkg/comment},
% \item \textsl{doclicense}, \url{https://ctan.org/pkg/doclicense},
% \item \textsl{draftwatermark}, \url{https://ctan.org/pkg/draftwatermark},
% \item \textsl{environ}, \url{https://ctan.org/pkg/environ},
% \item \textsl{etoolbox}, \url{https://ctan.org/pkg/etoolbox},
% \item \textsl{fancyhdr}, \url{https://ctan.org/pkg/fancyhdr},
% \item \textsl{float}, \url{https://ctan.org/pkg/float},
% \item \textsl{fontenc}, \url{https://ctan.org/pkg/fontenc},
% \item \textsl{framed}, \url{https://ctan.org/pkg/framed},
% \item \textsl{geometry}, \url{https://ctan.org/pkg/geometry},
% \item \textsl{graphicx}, \url{https://ctan.org/pkg/graphicx},
% \item \textsl{hyperref}, \url{https://ctan.org/pkg/hyperref},
% \item \textsl{hyperxmp}, \url{https://ctan.org/pkg/hyperxmp},
% \item \textsl{iftex}, \url{https://ctan.org/pkg/iftex},
% \item \textsl{libertine}, \url{https://ctan.org/pkg/libertine},
% \item \textsl{manyfoot}, \url{https://ctan.org/pkg/manyfoot},
% \item \textsl{microtype}, \url{https://ctan.org/pkg/microtype},
% \item \textsl{natbib}, \url{https://ctan.org/pkg/natbib},
% \item \textsl{newtxmath}, \url{https://ctan.org/pkg/newtxmath},
% \item \textsl{pbalance}, \url{https://ctan.org/pkg/pbalance},
% \item \textsl{refcount}, \url{https://ctan.org/pkg/refcount},
% \item \textsl{setspace}, \url{https://ctan.org/pkg/setspace},
% \item \textsl{totpages}, \url{https://ctan.org/pkg/totpages},
% \item \textsl{unicode-math}, \url{https://ctan.org/pkg/unicode-math},
% \item \textsl{xcolor}, \url{https://ctan.org/pkg/xcolor},
% \item \textsl{xkeyval}, \url{https://ctan.org/pkg/xkeyval},
% \item \textsl{xstring}, \url{https://ctan.org/pkg/xstring},
% \item \textsl{zi4}, \url{https://ctan.org/pkg/zi4},
% \item \textsl{zref-savepos}, \url{https://ctan.org/pkg/zref-savepos},
% \item \textsl{zref-user}, \url{https://ctan.org/pkg/zref-user}.
% \end{itemize}
%
%
%\subsection{Invocation and options}
%\label{sec:invocation}
%
% To use this class, put in the preamble of your document
% \begin{quote}
%   \cs{documentclass}\oarg{options}|{acmart}|
% \end{quote}
% There are several options corresponding to the type of the document and
% its general appearance.  They are described below.  Generally
% speaking, the options have |key=value| forms, for example,
% \begin{verbatim}
% \documentclass[format=acmsmall, screen=true, review=false]{acmart}
% \end{verbatim}
%
%
% The option |format| describes the format of the output.  There are
% several possible values for this option, for example,
% \begin{verbatim}
%   \documentclass[format=acmtog]{acmart}
% \end{verbatim}
% Actually the words |format=| can be omitted, e.g.,
% \begin{verbatim}
%   \documentclass[acmtog, review=false]{acmart}
% \end{verbatim}
% The possible formats are listed in Table~\ref{tab:opts_format}.
% Note that formats starting with |acm| are intended for journals,
% transactions, and course materials, while formats starting with
% |sig| are intended for proceedings published as books.
%
% Sometimes conference proceedings are published as a special issue
% (or issues) of an ACM journal.  In this case, you should use the
% journal format.  The templates |sample-acmsmall-conf.tex| and
% |sample-acmtog-conf.tex| show how to enter conference information.
% Note that you need to comment out |\acmJournal{...}| line for such
% papers to get the conference information in the footers and headers.
%
% \begin{table}
%   \centering
%   \caption{The possible values for the \texttt{format} option}
%   \label{tab:opts_format}
%   \begin{tabularx}{\textwidth}{>{\ttfamily}lX}
%     \toprule
%     \normalfont Value & Meaning\\
%     \midrule
%   manuscript & A manuscript. This is the default. \\
%       acmsmall & Small single-column format.  Used for ACMJCSS,
%            ACMJDS, AILET, CIE, CSUR,
%            FAC, GAMES, JACM, JATS,  JDIQ, JDS, JEA, JERIC,
%            JETC, JRC, PACMCGIT, PACMHCI, PACMMOD, PACMNET,
%            PACMPL, PACMSE, POMACS, TAAS, TACCESS, TACO, TAIS, TAISAP,
%            TALG, TALLIP (formerly TALIP), TCPS, TDS,
%            TEAC, TECS, TELO, THRI, TIIS, TIOT, TISSEC, TIST, TKDD, TMIS,
%            TOCE, TOCHI, TOCL,
%            TOCS, TOCT, TODAES, TODS, TOIS, TOIT, TOMACS, TOMM (formerly
%            TOMCCAP), TOMPECS, TOMS, TOPC, TOPLAS, TOPML, TOPS, TORS
%            TOS, TOSEM, TOSN, TQC, TRETS,
%            TSAS, TSC, TSLP, and TWEB, including special issues. \\
%   acmlarge  & Large single-column format.  Used for DLT, DTRAP, HEALTH,
%           IMWUT, JOCCH, and TAP, including special issues. \\
%   acmtog   & Large double-column format.  Used for
%          TOG, including annual conference Technical Papers.\\
%   sigconf & Proceedings format for most ACM
%          conferences (with the exception of SIGPLAN) and all ICPS
%          volumes.\\
%   sigplan & Proceedings format for SIGPLAN conferences.\\
%   acmengage & ACM EngageCSEdu Course materials.\\
%   acmcp & ACM cover page. \\
%   \bottomrule
%   \end{tabularx}
% \end{table}
%
% Starting in 2020, ACM retired formats |sigchi| and |sigchi-a|.
% SIGCHI conferences now use |sigconf| format for their publications.
% If a file uses |sigchi| format, a warning is issued, and the format
% is automatically switched to |sigconf|.  Format |sigchi-a| can be
% used for non-ACM documents only (see Section~\ref{sec:sigchi-a}).
% The format |acmcp| is used for ACM cover pages discussed in
% Section~\ref{sec:ug_acmcp}.
%
%
% There are several Boolean options that can take |true| or |false|
% values.  They are listed in Table~\ref{tab:opts_bool}.  The words
% |=true| can be omitted when setting a Boolean option, so instead of
% |screen=true| one can write just |screen|, for example,
% \begin{verbatim}
% \documentclass[acmsmall, screen, review]{acmart}
% \end{verbatim}
% The option |review| is useful when combined with the |manuscript| format
% option.  It provides a version suitable for reviewers and
% copy editors.
%
% Two samples in the |samples| directory, |manuscript| and
% |acmsmall-submission|, show manuscripts formatted for submission to
% ACM.
%
% The default for the option |screen| depends on the publication.  At
% present it is |false| for all publications \emph{but} PACM, since
% PACM is now electronic-only.  Thus PACM titles~(see
% Table~\ref{tab:pubs}) set this option to |true|.  In the future this
% option may involve additional features suitable for on-screen
% versions of articles.
%
% The option |natbib| is used when the corresponding
% \BibTeX\ style is based on |natbib|.  In most cases you do not need
% to set it.  See
% Section~\ref{sec:ug_bibliography}.
%
% The option |anonymous| is used
% for anonymous review processes and causes all author information to be
% obscured.
%
% The option |timestamp| is used to include a time stamp in the
% footer of each page.  When preparing a document, this can help avoid
% confusing different revisions.  The footer also includes the page range of
% the document.  This helps detect missing pages in hard copies.
%
% The option |authordraft| is intended for author's drafts that are not
% intended for distribution.  It typesets a copyright block to give the
% author an idea of its size and the overall size of the paper but
% overprints it with the phrase ``Unpublished working draft. Not for
% distribution.'', which is also used as a watermark.  This option sets
% |timestamp| and |review| to |true|, but these can be
% overriden by setting these options to |false| \emph{after}
% setting |authordraft| to |true|.
%
% The option |balance| determines whether the last page in the two
% column mode has balanced columns.  By default it is |true|; however,
% it may lead to problems for some documents.  When there are many
% figures near the end of the document, the attempts to balance
% columns may lead to the loss of the figures.  Set this option to |false| if
% you encounter problems.   An alternative is the
% (experimental) option |pbalance|, which uses the new package
% |pbalance|.  You may want to try |pbalance=true| to see if you get
% better results.  
%
% The option |urlbreakonhyphens| determines whether URLs can be split
% between lines after hyphens.  By default it is true.  Set it to
% |false| to disallow these breaks.
%
% \begin{table}
%   \centering
%   \caption{Boolean options}
%   \label{tab:opts_bool}
%   \begin{tabularx}{\textwidth}{>{\ttfamily}l>{\ttfamily}lX}
%     \toprule
%     \normalfont Option & \normalfont Default & Meaning\\
%     \midrule
%     review & false & A review version: lines are numbered and
%     hyperlinks are colored\\
%     screen & {\rmfamily see text} & A screen version:
%     hyperlinks are colored\\
%     natbib & true & Whether to use the |natbib| package (see
%     Section~\ref{sec:ug_bibliography})\\
%     anonymous & false & Whether to make author(s) anonymous\\
%     authorversion & false & Whether to generate a special
%     version for the authors' personal use or posting (see
%     Section~\ref{sec:ug_topmatter})\\
%     nonacm & false & Use the class typesetting options for
%     a non-ACM document, which will not include the conference/journal
%     header and footers.  Currenly such documents allow only a
%     Creative Commons license.\\
%     timestamp & false & Whether to put a time stamp in the
%     footer of each page\\
%     authordraft & false & Whether author's-draft mode is enabled\\
%     acmthm & true & Whether to define theorem-like environments, see
%                       Section~\ref{sec:ug_theorems}\\
%     balance & true & Whether to balance the last page in two column
%                        mode\\
%     pbalance & false & Whether to balance the last page in two column
%                        mode using pbalance package\\
%     urlbreakonhyphens & true & Whether to break urls on hyphens\\
%     \bottomrule
%   \end{tabularx}
% \end{table}
%
% The option |language| is used to define the languages for the
% multi-language papers.  It is discussed in
% Section~\ref{sec:ug_i13n}. 
%
%\subsection{Top matter}
%\label{sec:ug_topmatter}
%
% A number of commands set up \emph{top matter} or (in
% computer science jargon) \emph{metadata} for an article.  They
% establish the publication name, article title, authors, DOI and
% other data.  Some of these commands, like \cs{title} and \cs{author},
% should be put by the authors.  Others, like \cs{acmVolume} and
% \cs{acmDOI}---by the editors.  Below we describe these commands and
% mention who should issue them.  These macros should be used
% \emph{before} the \cs{maketitle} command.  Note that in previous
% versions of ACM classes some of these commands should be used before
% \cs{maketitle}, and some after it. Now they all must be used before
% \cs{maketitle}.
%
%
% This class internally loads the |amsart| class, so many top-matter
% commands are inherited from |amsart|~\cite{Downes04:amsart}.
%
% \DescribeMacro{\acmJournal}%
% The macro \cs{acmJournal}\marg{shortName} sets the name of the
% journal or transaction for journals and transactions.  The argument
% is the short name of the publication \emph{in uppercase}, for
% example,
% \begin{verbatim}
% \acmJournal{TOMS}
% \end{verbatim}
% The currently recognized journals are listed in
% Table~\ref{tab:pubs}.  Note that conference proceedings published in
% \emph{book} form do not set this macro.
%
%
%
% \DescribeMacro{\acmConference}%
% The macro
% \cs{acmConference}\oarg{short name}\marg{name}\marg{date}\marg{venue} is
% used for conference proceedings published in the book form.  The
% arguments are the following:
% \begin{description}
% \item[short name:] the abbreviated name of the conference (optional).
% \item[name:] the name of the conference.
% \item[date:] the date(s) of the conference.
% \item[venue:] the place of the conference.
% \end{description}
% Examples:
% \begin{verbatim}
% \acmConference[TD'15]{Technical Data Conference}{November
% 12--16}{Dallas, TX, USA}
% \acmConference{SA'15 Art Papers}{November 02--06, 2015}{Kobe, Japan}
% \end{verbatim}
%
% \DescribeMacro{\acmBooktitle}%
% By default we assume that conference proceedings are published
% in the book named \emph{Proceedings of \textsc{CONFERENCE}}, where
% \textsc{CONFERENCE} is the name of the conference infe